\section{Introduction}
A coordinate sweep of the Thorp shuffle acts on $N=2^d$ labeled cards.
It scans the $d$ binary coordinates in order and, in each direction,
independently swaps or leaves unchanged the two cards on every parallel edge.
A single card is uniform after a sweep. The full permutation retains
dependence because the cards have used the same switches. We prove that an
absolute number of sweeps suffices to mix the full permutation. Since a
sweep consists of $d$ physical shuffles, this gives the optimal order
$\Theta(\log N)$ in physical shuffle time.

The Thorp shuffle arose in a study of nonrandom shuffling and advantageous play in Faro \cite{Thorp1973}. For $N=2^d$, Morris proved an $O(d^{44})$ bound using evolving sets \cite{Morris2008}; Montenegro and Tetali sharpened the spectral-profile and evolving-set analysis to $O(d^{29})$ \cite[Section~6.4]{MontenegroTetali2006}. Morris then used entropy contraction to obtain $O((\log N)^4)$ mixing for every even deck size \cite{Morris2009}, followed by an $O(d^3)$ bound for dyadic decks \cite{Morris2013}. These estimates concern the full permutation and count individual physical shuffles; one coordinate sweep comprises $d$ such steps.

The $O(d^3)$ full-deck benchmark is also recorded in the November 2025 account of \citet[Section~1]{AlekseevEtAl2025}. Faster bounds for observing a fixed fraction of the cards, such as the $O((\log N)^2)$ bound of \citet{CzumajVocking2014}, concern a different projection of the permutation law. Here we obtain the order $O(d)$ for the full permutation on a dyadic deck, together with quantitative row--column representation estimates. The support bound in Proposition~\ref{ten:support} gives the matching lower order.

To see the geometric issue, split the coordinates into two consecutive
groups. The positions become a rectangular board. The first part of the
sweep acts within rows, and the second within columns; each part consists
of independent smaller sweeps. This suggests an induction on the number
of coordinates. Its obstacle is that the row and column symmetry
decompositions differ, and their projections generally do not commute.
We estimate how much of a vector selected by one decomposition can survive
selection by the other. Three forms of this estimate lead to the same
mixing conclusion while retaining different information about the
representation spaces.

\subsection{The sweep and its moments}
Positions are vectors $x=(x_1,\ldots,x_d)\in\mathbb F_2^d$. A physical
Thorp shuffle sends
\[
 (x_1,u)\longmapsto(u,x_1+\xi_u),
\]
where the $2^{d-1}$ bits $\xi_u$ are independent and fair. Let $q_d$
denote the law of this shuffle. Removing the
deterministic cyclic rotation after each shuffle gives successive matching
updates in the $d$ coordinate directions. After $d$ updates the rotation
is the identity. Thus one coordinate sweep has exactly the law of $d$
physical shuffles.

Permutations map initial positions to final positions, and $gh$ applies
$h$ first. Let $K_d=q_d^{*d}$ denote the law of one sweep. On the complex unitary
irreducible representation $V_\lambda$ of $S_N$, indexed by a partition
$\lambda\vdash N$, put
\[
 D_\lambda=\dim V_\lambda,\qquad
 K_d(\lambda)=\sum_{g\in S_N}K_d(g)\rho_\lambda(g).
\]
A matching layer is an orthogonal projection, so $K_d(\lambda)$ is a
contraction. All traces are ordinary, unnormalized traces, and
$\|A\|_p^p=\operatorname{Tr}|A|^p$. The matrix $K_d(\lambda)$ need not be
normal. On restriction to a subgroup, an irreducible type may occur
several times. We call its representation space the carrier and distinguish
directions within it from the different multiplicity copies.

Our first estimate attaches a positive weight to each Young diagram.
When we pass from the full board to its rows and columns, the difference
between the diagram weights pays for the cost of restoring deleted cells.
The weights still retain a fixed power of representation dimension.
Their construction appears in Section~\ref{ten:weighted:part:60}.

\begin{theorem}[Weighted sweep moments]\label{thm:weighted}
There are positive weights $W_\lambda$, an absolute finite $p_*$, and real
exponents $2\le p_d\le p_*$ such that, for every $d\ge1$ and every
$\lambda\vdash2^d$,
\begin{equation}\label{eq:weighted-main}
 D_\lambda^{3/4}\le W_\lambda\le D_\lambda,
 \qquad W_\lambda\|K_d(\lambda)\|_{p_d}^{p_d}\le1.
\end{equation}
Consequently $\|K_d(\lambda)\|_{\rm op}\le D_\lambda^{-3/(4p_*)}$.
\end{theorem}

The bounded exponent gives a dimension power uniformly over all diagrams
and deck sizes. It also yields a common unweighted moment: for an absolute
$q<\infty$,
\begin{equation}\label{eq:common-unit-moment}
 D_\lambda\|K_d(\lambda)\|_q^q\le1.
\end{equation}
Indeed, at $p=p_d$ the theorem gives
$\|K_d(\lambda)\|_{\rm op}^p\le W_\lambda^{-1}$, and hence
\[
 D_\lambda\|K_d(\lambda)\|_{4p/3}^{4p/3}
 \le D_\lambda W_\lambda^{-4/3}\le1.
\]
Since the matrices are contractions, $q=4p_*/3$ works simultaneously.
The later two routes prove a bound of the form
\eqref{eq:common-unit-moment} directly. Their additional information lies
in the intermediate row--column estimates described below.

We use
$\|\mu-\nu\|_{\rm TV}=\tfrac12\sum_g|\mu(g)-\nu(g)|$.
Let $t_{\rm mix}(d)$ be the least number of physical shuffles for which
the distance from the uniform full-deck law is at most $1/4$, from every
starting deck.
\begin{corollary}[Optimal order of mixing]\label{cor:optimal-mixing}
There is an absolute constant $C$ such that, for every $d\ge1$,
\[
 \left\lceil\frac2N\log_2\frac{3N!}{4}\right\rceil
 \le t_{\rm mix}(d)\le Cd,
 \qquad N=2^d.
\]
In particular $t_{\rm mix}(d)=\Theta(d)$, and its lower bound is
$2d-O(1)$.
\end{corollary}
The upper bound follows from finite-group Plancherel and Cauchy--Schwarz
\cite[Sections~2 and~3]{DiaconisShahshahani1981}, applied to a fixed
number of repeated forward sweeps. The lower bound counts the possible
permutations produced by the random switches. Both deductions are given
in Section~\ref{ten:weighted:part:285}, after the first moment proof.

\subsection{From row--column overlap to a bounded exponent}
The first route begins by choosing a subdiagram $a\subseteq\lambda$
inside an $(h,h)$-hook: no box of $a$ has both row and column index
greater than $h$. If $t=N-|a|$, we let the symmetric group on the
$N-t$ occupied cells of a board act through $V_a$. Hook trimming and the
choice of weights control the cost of making this reduction. The
representation $V_a$ occurs in a tensor space with an $h$-dimensional
even part and an $h$-dimensional odd part. This hook realization is
classical in signed tensor representation theory \cite{BereleRegev1987};
we obtain the form needed here from ordinary Schur--Weyl duality on the
two parts and signed induction.

Theorem~\ref{thm:occupied} bounds overlap on every occupied-cell pattern,
summing over all orthogonal copies of a fixed row type and column type.
Its proof majorizes the subgroup projections by positive mixtures of
products of one-cell states. Normalizing by the average state on each side reduces
the product estimate to the entropy of a uniformly chosen occupied cell
relative to the product of its row and column marginals. This is where
the geometry of the board enters: the entropy cost is controlled by the
number of deleted cells, regardless of their arrangement. The entropy step
is a finite form of the duality between product inequalities and entropy
subadditivity \cite[Theorem~2.1 and (2.4)]{CarlenCorderoErausquin2009}.

To return to $V_\lambda$, choose a diagram $b$ on the $t$ deleted boxes
such that $V_\lambda$ occurs in the representation induced from
$V_a\otimes V_b$. That induced representation has one copy of
$V_a\otimes V_b$ for each placement of the deleted cells. Restricting a row or column representation to the stabilizer of a
placement decomposes its carrier into tensor products of surviving and
deleted carrier spaces, with multiplicities. The resulting coefficient states can be entangled between these
carrier factors. Lemma~\ref{lem:entangled} transfers the occupied-cell
bound to such states. Positive coefficient states have total trace one, so summing
their contributions introduces no new carrier-dimension count. The
remaining placement entropy is paid by the difference between the
weight of $\lambda$ and the weights of its row and column types.

\begin{figure}[ht]
\centering
\begin{tikzpicture}[scale=.62]
\foreach \x in {0,...,6} {\draw[gray!60] (\x,0)--(\x,4);}
\foreach \y in {0,...,4} {\draw[gray!60] (0,\y)--(6,\y);}
\foreach \x/\y in {1/0,4/1,2/2,5/3} {\fill[gray!35] (\x,\y) rectangle ++(1,1);\draw (\x+.2,\y+.2)--(\x+.8,\y+.8);\draw (\x+.2,\y+.8)--(\x+.8,\y+.2);}
\draw[-{Stealth},thick,blue] (-.8,1.5)--(6.8,1.5);
\draw[-{Stealth},thick,red] (3.5,-.7)--(3.5,4.7);
\node[blue] at (7.4,1.5) {rows};\node[red] at (3.5,5.1) {columns};
\node[align=left,anchor=west] at (8,3.1) {occupied cells: $V_a$\\deleted cells: $V_b$\\one placement: $V_a\otimes V_b$};
\end{tikzpicture}
\caption{A board with deleted cells, marked by crosses. The occupied-cell estimate applies to every deletion pattern. Restoring the deleted cells assigns the representation $V_a\otimes V_b$ to each placement; the weight comparison pays for summing over placements. The arrows indicate row and column subgroup actions.}
\label{fig:board}
\end{figure}

Weighted Schatten interpolation then transfers the projection bound
to arbitrary matrices in the row and column Fourier blocks. Its
logarithmic error is independent of the Schatten exponent. Applied to
the two families of smaller sweeps, the bound increases the required
exponent by only a factor $1+C/d$. These factors have bounded product
along successive halvings of $d$. A direct forest estimate covers
diagrams whose first row contains almost every box; a strict norm gap
in finitely many small dimensions starts the induction. This proves
Theorem~\ref{thm:weighted} in Section~\ref{sec:completion}.

\subsection{Two further forms of the overlap estimate}
Section~\ref{ten:frames} keeps track of the number of selected directions
inside each row and column carrier, while including every multiplicity
copy. Its Schatten-$8$ bound has a cost controlled by the number of boxes
below the first row. A separate signed-tensor estimate has a cost
controlled by the number of boxes outside a small hook. The two bounds
cover complementary shapes, and a sparse estimate using short time
windows completes another bounded-exponent induction.

Section~\ref{rec:sec-coherent} instead retains the deleted positions as
an ordered list. Summing squared coefficients over the initial,
intermediate, and final hole lists gives the explicit restoration
factor $e^{3l}\binom Nl$, where $l$ is the number of holes. Its
occupied-cell estimate resolves the row and column projections by averaging
over compact-group orbits of highest-weight vectors
\cite[Section~2]{Perelomov1972}. A minimum-norm argument matches their
averaged one-cell densities to a common density before the row--column
comparison. This normalization is related to the highest-weight capacity
framework of \citet[Sections~2 and~4]{FranksWalter2022}; we prove the
required form locally to keep the loss from missing cells explicit. The
overlap proof and moment induction are given here.
The sparse part uses the partial-deck density and lower-diagram bounds
of \citet[Theorem~6.1 and Proposition~6.3]{OpenAI2026Routing}.

The companion \emph{Random-subspace tests and trace smoothing for
coordinate sweeps} \cite{OpenAI2026RandomSubspace} develops related
trace-smoothing estimates. For symmetric-group and tableau conventions we
follow \citet{Sagan2001,Stanley1999}; branching and content spectra are
also developed by \citet{VershikOkounkov2005}. We use ordinary
Schur--Weyl duality in the form of \citet[Theorem~4.57 and
Corollary~4.59]{etingof}.

Section~\ref{sec:conventions} fixes the common representation conventions.
Part~\ref{part:weighted} constructs the weights, proves the occupied-board
and restoration bounds, and completes the first moment proof and the
mixing deduction. Part~\ref{part:frames} develops the selected-rank frame
estimates and their moment induction. Part~\ref{part:coherent} proves
the ordered-hole estimate and closes the coherent-state induction using
the cited partial-deck input.
